\documentclass[12pt,a4paper]{article}
% Known tension (kept deliberately): because the inequality base I and the convergence base C are
% built independently (I = current distribution rescaled, then G = I after a global tax; C = Bothe's
% converged income, then N = C minus the global transfer rescaled to the no-global GDP level), and
% because the no-global GDP exceeds the converged GDP, we do NOT have the accounting identity
% C - N = G - I. Occasionally this makes G slightly larger than C (or N larger than I). We accept it:
% the national redistribution that turns G into C also lowers GDP, so a few people who gained from
% I->N need not gain again from G->C. Forcing G = I + C - N would risk negative incomes at the bottom
% in 2100.
% we don't have C - N = G - I, hence sometimes illogical things happen (e.g. B90k IT, 2500-2800€/m), s.t. G > C and N > I (> G). This is because I = current rescaled (then G = I⋅(1-t(I))) are defined separately from C (then N = (C - GR)⋅GDP_I): since GDP I~G > GDP C~N, we can have G > C despite N < I. We can justify it by saying that the national redistr occurring after G reduces GDP further, so some who benefited from I->N don't benefit from G->C. 
% We cannot set G = I + C - N because in 2100 IT p0 it could yield 0 = 1k + 2k - 3k. In other words, N couldn't be that large if we also have to have G; or C embeds as lower GDP cumulative effects of G which go beyond the GDP reduction due to N. 
\usepackage[T1]{fontenc}
\usepackage[utf8]{inputenc}
\usepackage[english]{babel}
\usepackage{lmodern}
\usepackage{microtype}
\usepackage{geometry}
\geometry{margin=2.5cm}
\usepackage{amsmath,amssymb}
\usepackage{booktabs}
\usepackage{array}
\usepackage{multirow}
\usepackage{tabularx}
\usepackage{xcolor}
\usepackage{hyperref}
\hypersetup{colorlinks=true, linkcolor=blue, citecolor=blue, urlcolor=blue}
\usepackage[round,sort&compress]{natbib}
\usepackage{setspace}
\onehalfspacing

% ─── Convenience macros ───────────────────────────────────────────────────────
\newcommand{\note}[1]{%
  \begin{center}
    \colorbox{yellow!30}{\parbox{0.9\linewidth}{\textbf{Note:} #1}}
  \end{center}}
\newcommand{\scenario}[1]{\texttt{#1}}
\newcommand{\col}[1]{\texttt{#1}}
\newcommand{\file}[1]{\texttt{#1}}
\newcommand{\EUR}{\texteuro}

% ─── Document ─────────────────────────────────────────────────────────────────
\begin{document}

\title{Appendix: Methodology for the Conjoint Analysis Question}
\author{Adrien Fabre \and Viola Maone}
\date{\today}
\maketitle
\tableofcontents
\bigskip

% ==============================================================================
\section{Overview}
\label{sec:overview}
% ==============================================================================

The conjoint questions ask respondents to choose between two scenarios for the century. Every scenario determined by a few parameters: working hours, whether there is global and/or national redistribution, the pace of decarbonization, and structural change. Each scenario is presented to the respondent with concrete outcomes: their own projected \emph{cash income in 2035}, the \emph{income distributions in 2100} for Italy and for the world, and the \emph{global temperature in 2100}.
This appendix describes how those three outcomes are computed. 
% Almost all the work is done once, offline, by a single R script, \file{code\_simulator/questionnaire.R}. The script turns the raw
% research data of \citet{Bothe2026} and \citet{Chancel2026} into a handful of small CSV files: one
% table of 2035 cash-income distributions, one table of 2100 income distributions, one table of 2100
% temperatures, and one table of scalar constants. The web page then merely \emph{reads} these tables
% and, at display time, interpolates and rescales them for the particular scenario shown --- no
% economics is recomputed in the browser. Reading this appendix is therefore equivalent to reading
% \file{questionnaire.R}: the sections below follow the script's own logic, from the raw inputs to the
% exported files.

% All monetary amounts are expressed in EUR at 2025 purchasing power, per adult, per year, unless stated otherwise.

% ==============================================================================
\section{The three conjoint tasks}
\label{sec:tasks}
% ==============================================================================

The Italian survey runs three variants of the conjoint: %, all living in
% \file{code\_simulator/IT\_survey/}. 
% They share the same engine
% (\file{IT\_survey/scenarios.js}) and the same precomputed data; they differ only in how the two
% scenarios on screen are chosen and in which 2035 dataset they read.

\begin{description}
  \item[\file{conjoint\_few}.] The respondent compares two scenarios drawn from a list of twelve variants of the Baseline Ethical Convergence scenario (or simply BC), where we depart from BC either in terms of working hours (BC45k, BC90k, BC120k), decarbonization pace (BC\_SD), redistribution scope (BI, BN, BG), beef (BCbeef -- no decrease in beef consumption), flight consumption (BCflight), or the evolution of public services (BCmat -- no sectoral recomposition). We also include the Benchmark Inequality scenario, with increased working hours, no redistribution and no behavioral nor structural change (BI120k). The default BC scenario is drawn twice as often as the others.

  \item[\file{conjoint\_any}.] This explores the full design space rather than a curated shortlist: the parameters of each scenario are drawn uniformly and independently (the draw is repeated if the two sides come out identical). 

  \item[\file{conjoint\_any2}.] Similar to \file{conjoint\_any} except that hourly productivity grows at Italy's  2025--2035 productivity rate modelled by \cite{Chancel2026} (close to 3\%), instead of the flat 1.5\,\%/year assumption used by \file{conjoint\_few} and \file{conjoint\_any}.
\end{description}

% Because \file{conjoint\_few} and \file{conjoint\_any} share the flat-growth dataset and \file{conjoint\_any2} uses the country-growth dataset, the script exports both: \file{ineq\_IT\_2035} and \file{ineq2\_IT\_2035}. Everything else --- the 2100 distributions, temperatures and constants --- is common to all three tasks.

% ==============================================================================
\section{How scenarios are named}
\label{sec:naming}
% ==============================================================================

A scenario combines four structural choices.

\begin{description}
  \item[Redistribution scope.] Which redistribution operates: \scenario{C} = Convergence (both national and global), \scenario{G} = Global only, \scenario{N} = National only, \scenario{I} = Inequality (neither). %By convention the scope letter is folded into the class label, e.g.\ \scenario{B90kC}, \scenario{SN}, \scenario{PI}.
  \item[Working hours (the ``class'').] The weekly hours \emph{worked in 2035} --- 32, 36, 40 or 44 --- which the model carries forward to a 2100 GDP per capita 45, 60, 90 or 120 thousand euros in Convergence scenarios, respectively (Table~\ref{tab:hours}). 36\,h is the  benchmark; the longer hours classes are labelled \scenario{B90k} (40\,h) and \scenario{B120k} (44\,h, corresponding to PC/PI scenarios in \citealp{Chancel2026}), the shorter \scenario{B45k} (32\,h). %The data files also carry a still-shorter 28\,h class, unused by the Italian survey.
  \item[Sectoral change.] \scenario{1} = no change (public services stable, current diet);
    \scenario{2} = structural change (expanded public services, less beef and fewer flights).
  \item[Decarbonisation.] \scenario{SD} = slow, \scenario{ID} = intermediate, \scenario{FD} = fast, matching the IEA CPS / STEPS / NZE pathways.
\end{description}

\noindent So \scenario{BC2\_FD} is the default benchmark (36\,h worked, full convergence, structural change, fast decarbonisation), and \scenario{BI90k1\_FD} is the 40\,h, no-redistribution, no-structural-change, fast-decarbonisation scenario.

\begin{table}[ht]
\centering
\caption{Working hours, class label, and the resulting 2100 GDP per capita.}
\label{tab:hours}
\begin{tabular}{lccc}
\toprule
\multirow{2}{*}{Hours worked, 2035} & \multirow{2}{*}{Class label}
  & \multicolumn{2}{c}{World 2100 GDP per capita (k\EUR)} \\
\cmidrule(lr){3-4}
 & & With global redistr. & Without \\
\midrule
32 & \scenario{B45k}            & 45  & 37.5 \\
36 & \scenario{B} (\scenario{S}) & 60 & 50   \\
40 & \scenario{B90k}            & 90  & 75   \\
44 & \scenario{B120k} (\scenario{P}) & 120 & 100 \\
\bottomrule
\end{tabular}
\end{table}

% \note{%
%   \citet{Bothe2026} provide the income distribution only in the Convergence scenario. The variants that differ only in decarbonisation pace (\scenario{SD}/\scenario{ID}/\scenario{FD}) or in sectoral change (\scenario{1}/\scenario{2}) all start from the same income distribution; we layer the differences on top ourselves, through an income-cost factor for decarbonisation (Section~\ref{sec:factors}) and a flat-rate tax for expanded public services
%   (Section~\ref{sec:factors}).
% }

% ==============================================================================
\section{Raw data sources}
\label{sec:data}
% ==============================================================================

% The script reads 
We use 
five raw input datasets and one pre-computed helper file:

\begin{description}
  \item[\citet{Bothe2026} simulation.] For every country, income percentile and year, we have income after national taxes and transfer, the global income tax paid, as well as national net income per adult and the uniform global dividend received. From these we reconstruct, for each cell, the \emph{full income} under global convergence:
    \[
      \text{full income} \;=\; \text{income after national taxes and transfers}
      \;-\; \text{global income tax} \;+\; \text{global dividend received}.
    \]
    The last two terms make up the \emph{net global redistribution} (positive for  poor people, negative for rich ones). %Removing them recovers the pre-global ``national'' income; this is the lever that lets us switch global redistribution on and off  (Section~\ref{sec:scopes}).
  \item[\citet{Chancel2026} macro scenarios.] Per-country, per-year time series for GDP (total and  per capita, under both the convergence and the inequality paths), labour productivity, hours  worked per capita and per worker, public-spending shares, and projected populations. These supply the hours/productivity scalings, the public-services tax, and the populations used to  pool countries into a world distribution.
  \item[\citet{FisherGethin2023} for Italy.] The pre-tax income shares and the various taxes and transfers by percentile, from which we build the \emph{gross cash income} of Italians in 2025 (Section~\ref{sec:cash2025}). 
  \item[WID retention ratio and \citet{Garbinti2021} wealth composition.] Two small ingredients of  the cash-income calculation: the share of corporate profits retained rather than paid out, and  how housing wealth is spread across the distribution (for imputed rents).
  \item[Pre-computed Chancel temperatures.] A table of 2100 temperatures for the scenarios \citet{Chancel2026} explicitly model, used directly and as the training set for the few scenarios they do not model (Section~\ref{sec:temperature}).
\end{description}

% ==============================================================================
\section{Italian cash income in 2025: the anchor}
\label{sec:cash2025}
% ==============================================================================

Respondents think in terms of the money that actually reaches their household, so every distribution shown to them is expressed in \emph{gross cash income}: income at market prices, net of all taxes (including VAT), imputed rents and retained earnings, but before deducting the consumption of fixed capital. %The whole pipeline is anchored to one observable quantity --- the 2025 cash income of Italians by percentile --- so that the figure a respondent recognises as ``their income'' lines up with the model.

Starting from the \citet{FisherGethin2023} pre-tax income share at each percentile, the script subtracts direct taxes (income, wealth and corporate), social contributions and indirect taxes (VAT), adds back cash transfers from government, then makes three small adjustments: it removes the imputed rent on owner-occupied housing (using the \citealt{Garbinti2021} housing-wealth profile), removes the retained (non-distributed) part of corporate profits, and adds back capital consumption for the capital-income share. %Finally, it subtracts the global tax already paid in 2025. 
The result is scaled to euros by Italy's 2025 national net income per adult from \cite{Bothe2026}. The bottom 99\,\% of the resulting profile is sorted into increasing order to remove imputation noise (the top percentiles are left untouched).

Concretely, with $\bar a=\sum_p w_p\,a^{\mathrm{pre}}_p$ the population-mean pre-tax income and $\kappa_p$ percentile $p$'s share of capital income, the 2025 cash income is
\[
  c_{25}(p)=\mathrm{NNI}\left[\frac{a^{\mathrm{pre}}_p-T_p+S_p}{\bar a}-\rho\,\kappa_p-y\,h_p\right]+\kappa_p\,\mathrm{CFC},
\]
where $T_p$ gathers direct, corporate, social and indirect (VAT) taxes, $S_p$ government cash transfers, $\rho$ the retained-earnings rate, $y$ the imputed-rent yield, $h_p$ the housing-wealth share, $\mathrm{NNI}$ Italy's 2025 net national income per adult and $\mathrm{CFC}$ the per-adult consumption of fixed capital.

\note{%
  VAT is subtracted because the \citet{FisherGethin2023} pre-tax income is measured at market prices, which include indirect taxes. Removing VAT converts it to factor-cost income, the basis on which cash incomes add up. Skipping this step would overstate cash income.
}

\note{%
  Why anchor everything to 2025 cash income? Because it is the only quantity in the whole exercise that respondents can verify against their own pay slip. Every later distribution is built so that, at the percentiles, it remains comparable to this anchor.
}

% ==============================================================================
\section{Income distributions in 2035}
\label{sec:it2035}
% ==============================================================================

\subsection{One base, four multiplicative factors}
\label{sec:factors}

All 2035 distributions descend from a single base, \col{IT35}: the gross cash income Italians would have in 2035 under the convergence benchmark with slow decarbonisation and no structural change (\scenario{BC1\_SD}). We obtain it by taking \citet{Bothe2026}'s 2035 SC income distribution for Italy and rescaling each percentile by that percentile's 2025 cash-to-income ratio. % --- so the \emph{shape} comes from the simulation while the \emph{level} stays anchored to observable 2025 cash income.

Every other 2035 scenario is then this base (or a scope-specific cousin of it, see below) multiplied by up to four independent factors:

\begin{description}
  \item[An hours factor.] Hours proportionately translate into output. The 40\,h ``\scenario{B90k}'' class is the family's reference (hours are held at their 2025 level); the 32\,h and 44\,h classes are \scenario{B90k} scaled by their worked-hours ratio (32/40, 44/40), and the 36\,h ``\scenario{B}'' class follows \cite{Chancel2026}'s modelled hours trajectory (close to $36/40$), reaching exactly the 60k income level by 2100. %Income is thus monotone in hours within the family. \scenario{B90k}'s 40\,h income still lands slightly \emph{below} the \scenario{SC} benchmark (Section~\ref{sec:bfamily}), but only because \scenario{SC} assumes \citeauthor{Chancel2026}'s faster 3\,\%/year productivity growth, not because of hours.
  \item[A scope adjustment.] Turning national and/or global redistribution on or off both reshapes the distribution and moves its average to the GDP level that scope implies (Section~\ref{sec:scopes}).
  \item[A decarbonisation factor.] A flat income penalty for faster decarbonisation: 1.00 for slow, 0.99 for intermediate, 0.97 for fast --- i.e.\ fast decarbonisation costs 3\,\% of income at each percentile in 2035. We choose these factors to reflect the extra investment diverted into the energy transition, noting that the intermediate scenario is closer to the slow pace than to the fast one. 
  \item[A public-services factor.] When public services expand (the structural-change option), the extra spending is financed by a flat-rate tax. The rate is set so that the additional public-spending share of national expenditure between 2025 and 2035 is exactly covered by taxing average income at that rate; for Italy it is about 4.9\,\%. The factor is then one minus this rate.
\end{description}

\noindent In symbols, the base is the 2035 SC income carried into cash units percentile by percentile, and each convergence-scope scenario at hours class $k$ multiplies it by the three factors:
\[
  \col{IT35}(p)=y^{\mathrm{SC}}_{2035}(p)\,\frac{c_{25}(p)}{y_{2025}(p)},
  \qquad
  \mathrm{col}_k(p)=\col{IT35}(p)\;\eta_k\;\delta\;\psi,
\]
with the per-percentile cash-to-income ratio $c_{25}/y_{2025}$ (here $y_{2025},y^{\mathrm{SC}}_{2035}$ are \citeauthor{Bothe2026}'s 2025 and 2035 SC incomes), hours factor $\eta_k$ (Section~\ref{sec:bfamily}), decarbonisation factor $\delta\in\{1,0.99,0.97\}$, and public-services factor $\psi=1-\tau$ ($\tau\approx4.9\%$) when services expand, $\psi=1$ otherwise.

\noindent The 36\,h benchmark \scenario{B} is \col{IT35} with its hours factor and the fast-decarbonisation and public-services factors applied; \scenario{B90k} (40\,h), \scenario{B45k} (32\,h) and \scenario{B120k} (44\,h) only change the hours factor, while the \scenario{N}, \scenario{G} and \scenario{I} scope variants substitute the appropriate base and scope adjustment. The no-redistribution (inequality-scope) scenarios are the exception: \scenario{BI} (36\,h), \scenario{BI120k} (44\,h) and their hours variants start from 2025 cash income rather than from \col{IT35}, because they keep the current distribution instead of the converged one --- the 2025 cash income is rescaled to the no-global GDP level and to each hours class, after which the same decarbonisation and public-services factors apply.

% \subsection{The net global transfer, GR35}
% \label{sec:gr35}

% Alongside the income columns, the script stores \col{cash\_GR35}: the per-percentile \emph{net global transfer} in 2035 (dividend received minus global tax paid). %, expressed in the same cash units. 
% % It is positive at the bottom of the distribution (net receivers) and negative at the top (net contributors). 
% % This column is the explicit ingredient that builds the national-only and global-only scopes, so --- unlike the income columns --- it is deliberately \emph{not} sorted into increasing
% % order: it is meant to be decreasing in income.

\subsection{The four redistribution scopes}
\label{sec:scopes}

The four scopes are best understood as the four cells of a $2\times2$ table (national redistribution on/off $\times$ global redistribution on/off), all sitting at one of two average-income levels: the converged level (when global redistribution is on) or the higher ``no-global'' level (when it is off, since the world economy converges less and stays richer on average).

% Alongside the income columns, the script stores \col{cash\_GR35}: the per-percentile \emph{net global transfer} in 2035 (dividend received minus global tax paid). %, expressed in the same cash units. 

\begin{description}
  \item[\scenario{C} (Convergence).] This uses the \citet{Bothe2026} full income directly: \col{IT35} and its hours/decarbonisation/public-services variants.
  \item[\scenario{I} (Inequality).] It uses the current within-country shape, rescaled to the no-global GDP level. 
  % This is the base from which the other ``no-redistribution-shape''     scenarios are grown.
  \item[\scenario{N} (National only).] It takes the   nationally-redistributed income (the convergence income minus the net global redistribution --- dividend received minus global tax paid) and rescales it to the no-global GDP level. It has the compressed, redistributed shape but the higher no-global average.
  \item[\scenario{G} (Global only).] It keeps the inequality shape, then apply a \emph{global} transfer schedule on top: at each income level the schedule subsidises the poor and taxes the rich at the rate the global transfer does worldwide (to move from \scenario{N} to \scenario{C}). That schedule is computed at the world level but applied to each Italian income.
\end{description}

Writing $\mathrm{GDP}_X=\sum_p w_p\,X(p)$ for a base's population average, $\scenario{I}_0=c_{25}\,s$ for the 2025 cash income rescaled to the no-global GDP level (factor $s$) and $\mathrm{GR}$ for the net global transfer \col{cash\_GR35}, the four 2035 bases (before the $\delta,\psi$ factors) are
\[
  \scenario{C}=\col{IT35},\quad
  \scenario{I}=\scenario{I}_0,\quad
  \scenario{N}=(\col{IT35}-\mathrm{GR})\,\frac{\mathrm{GDP}_{\scenario{I}_0}}{\mathrm{GDP}_{\col{IT35}}},\quad
  \scenario{G}=(\scenario{I}_0+\mathrm{GR})\,\frac{\mathrm{GDP}_{\col{IT35}}}{\mathrm{GDP}_{\scenario{I}_0}}.
\]

\note{%
  The global schedule for \scenario{G} is read off the \emph{world} distribution (what rate each income level faces globally) and then applied within Italy. As a result the neat accounting identity ``\scenario{G} is to \scenario{I} as \scenario{C} is to \scenario{N}'' holds at the world level but not within a single country, and occasionally we can have both \scenario{G}$>$\scenario{C} and \scenario{N}$>$\scenario{I}. This can be interpreted as a further reduction of GDP when national redistribution occurs after global redistribution, so that some who benefit from national redistribution in absence of global redistribution don't benefit from it when global redistribution is present. 

  We keep this rather than forcing the identity $\scenario{C} - \scenario{G} = \scenario{N} - \scenario{I}$, which could otherwise drive the lowest 2100 incomes negative in \scenario{G}. Indeed, \scenario{C} and \scenario{G} embed as a lower GDP the cumulative effects of global redistribution, leaving less space to increase the lowest income from \scenario{G} to \scenario{G} than with from \scenario{I} to \scenario{N}. 
  
}
% G > C and N > I (> G). This is because I = current rescaled (then G = I⋅(1-t(I))) are defined separately from C (then N = (C - GR)⋅GDP_I): since GDP I~G > GDP C~N, we can have G > C despite N < I. We can justify it by saying that the national redistr occurring after G reduces GDP further, so some who benefited from I->N don't benefit from G->C. 
% We cannot set G = I + C - N because in 2100 IT p0 it could yield 0 = 1k + 2k - 3k. In other words, N couldn't be that large if we also have to have G; or C embeds as lower GDP cumulative effects of G which go beyond the GDP reduction due to N. 

\subsection{The B family} % and why \texorpdfstring{\scenario{B}}{B} differs from \texorpdfstring{\scenario{B90k}}{B90k}}
\label{sec:bfamily}

Compared to \cite{Chancel2026} scenarios which entail a 3\% hourly productivity growth over 2025--2035 for Italy, the ``B'' scenarios use a more conservative 1.5\% growth rate. Scenario \scenario{B90k} (40\,h) is the family's reference. It holds hours per worker at their 2025 level while sharing Convergence's long-run productivity path. 
% , and its 2035 income is the benchmark scaled by a factor that nets out convergence's front-loaded productivity and its falling hours. 
The shorter and longer classes (\scenario{B45k} at 32\,h,
\scenario{B120k} at 44\,h) are \scenario{B90k} times their worked-hours ratio (e.g., $\scenario{B45k} = \scenario{B90k} \cdot 32/40$). The 36\,h class \scenario{B} is special: it uses the working hours trajectory of \cite{Chancel2026} (close to $36/40$ times 2025 hours per worker in 2035); by 2100 it reaches exactly the 60k income level from \cite{Chancel2026}.

With $r=1.5\%$ the B growth rate, $\gamma$ Italy's 2025--2035 SC productivity growth and $\Delta h$ SC per-worker hours change over the same period, the hours factors of Section~\ref{sec:factors} are
\[
  \eta_{\scenario{B90k}}=\frac{(1+r)^{10}}{\Delta h\,\gamma},\qquad
  \eta_{\scenario{B}}=\Delta h\,\eta_{\scenario{B90k}},\qquad
  \eta_{\scenario{B45k}}=\tfrac{32}{40}\,\eta_{\scenario{B90k}},\qquad
  \eta_{\scenario{B120k}}=\tfrac{44}{40}\,\eta_{\scenario{B90k}}.
\]
The $\Delta h\,\gamma$ in the denominator removes convergence's front-loaded productivity and its falling hours, leaving \scenario{B90k} on the conservative flat-growth path.

% \note{%
%   Because convergence productivity growth peaks mid-century and then declines, the 2035 productivity
%   gain can exceed the flat 1.5\,\%/year assumed for the B scenarios. When it does, \scenario{B90k}
%   income in 2035 comes out \emph{below} the benchmark --- a sensible outcome, since the benchmark
%   squeezes peak productivity into a 36\,h week while \scenario{B90k} works 40\,h at a lower
%   per-hour frontier.
% }

% ==============================================================================
\section{Income distributions in 2100}
\label{sec:2100}
% ==============================================================================

\subsection{Full income, rescaled by one number}

The meaningful comparison across scenarios and countries is ``full income'' (secondary income, inclusive of in-kind public services and collective expenditures), not cash income. To keep the 2100 charts on the same footing as the 2035 cash income the respondent has just considered, every 2100 distribution is multiplied by a scalar: the ratio of average cash income to average full income for Italy in 2025, which is about 0.69. This one scalar --- not the per-percentile ratio used in 2035 --- is applied to every percentile and every country, so the 2100 distributions reflect genuine inequality rather than Italy's percentile-by-percentile cash-to-income quirks.

\note{%
  Using a per-percentile ratio in 2035 (to match the respondent's own income) but a single scalar in 2100 (to keep cross-country comparisons clean) is a deliberate choice. The 2035 figure must line up with one person's pay slip; the 2100 figure must let Italy and the world be compared on the same basis.
}

\subsection{Italy and the world}

For Italy, the four scopes at 2100 are built exactly as in 2035, from \citet{Bothe2026}'s 2100 income: Convergence is the full income itself; National-only rescales to the no-global levels (from \citet{Bothe2026}'s Persistent Inequality scenario); Inequality grows the 2025 income by the inequality path's GDP factor and halves it (the no-redistribution level); Global-only applies the world transfer schedule to the Inequality income.

With $Y_{2100}$ \citeauthor{Bothe2026}'s SC income, $D$ the global dividend, $y^{\mathrm{nat}}_{2025}$ the pre-global 2025 national income, $\Gamma_{\scenario{I}}$ the inequality path's 2025--2100 GDP growth and $\tau_G(\cdot)$ the world tax-rate schedule (a function of national income, read off the world \scenario{SN}$\to$\scenario{SC} distributions),
\[
  \scenario{C}=Y_{2100},\quad
  \scenario{N}=(Y_{2100}-D)\,\frac{\mathrm{GDP}_{\scenario{I}}}{\mathrm{GDP}_{\scenario{C}}},\quad
  \scenario{I}=\tfrac12\,y^{\mathrm{nat}}_{2025}\,\Gamma_{\scenario{I}},\quad
  \scenario{G}=\bigl(1-\tau_G(\scenario{I})\bigr)\scenario{I}.
\]
Each column is then multiplied by the cash-equivalent scalar ($\approx0.69$) and, for the hours classes, by the Table~\ref{tab:coefs2100} multiplier.

To aggregate countries' distributions into the world's one, we pool every (country, percentile) cell weighted by its population share, sort them by income, and read off one hundred equal-population global percentiles (filling any gaps by interpolation). The population used follows the scenario: the Convergence population (about 9.4 billion in 2100) when global redistribution is on, the Inequality population (about 10.2 billion) when it is off.

The hours classes do not need a separate 2100 calculation: each maps to a fixed 2100 GDP target, so its distribution is simply the scope's distribution scaled by target over 60 (Table~\ref{tab:coefs2100}). %The 36\,h class \scenario{B} --- and its beef/flight/material variants --- equals convergence at 2100.

\begin{table}[ht]
\centering
\caption{2100 income multipliers relative to convergence (\scenario{SC}), applied to both Italy and
  the world. \scenario{B120k} and \scenario{P} share the multiplier 2.0; they differ only in their
  2035 path. The food variants of \scenario{B} equal \scenario{SC} at 2100.}
\label{tab:coefs2100}
\begin{tabular}{lcccc}
\toprule
Class & \scenario{B45k} & \scenario{S}/\scenario{B} & \scenario{B90k} & \scenario{B120k},\,\scenario{P} \\
Hours worked, 2035 & 32 & 36 & 40 & 44 \\
2100 GDP target (k\EUR) & 45 & 60 & 90 & 120 \\
\midrule
Multiplier & 0.75 & 1.00 & 1.50 & 2.00 \\
\bottomrule
\end{tabular}
\end{table}

% \note{%
%   The inequality (\scenario{PI}) base for 2100 uses the pre-global 2025 income, not the income that already includes the 2025 global dividend. Using the post-dividend figure would inflate the inequality distribution for net receivers and depress it for net contributors --- exactly the redistribution the inequality scenario is supposed to lack.
% }

\note{%
  Structural change leaves 2100 income unchanged: by 2100 the convergence economy reaches the same equilibrium income whether or not the food system and public services were reformed earlier. So \scenario{SC} and \scenario{SCmat} have identical 2100 distributions (structural change shows up in 2035 income, via the decarbonisation factor, and in temperature).
}

% ==============================================================================
\section{Temperature in 2100}
\label{sec:temperature}
% ==============================================================================

Temperature depends on the world 2100 GDP, on the decarbonisation pace, and on structural change. For the scenarios \citet{Chancel2026} explicitly model, the script uses their reported 2100 temperature directly. For the few scenario types they do not model (the no-global and 40\,h families), it predicts temperature from a regression of the observed temperatures on 2100 GDP, decarbonisation pace and structural change, fitted on the two dozen scenarios \citet{Chancel2026} do
report. The fit is very tight (it reproduces the observed temperatures to within 0.1\,°C, the rounding granularity). %, though its accuracy on the extrapolated scenarios is necessarily unknown. TODO! find a better fit?

The base prediction is the OLS fit
\[
  T = \beta_0 + \beta_g\,G + \sum_{d\in\{\mathrm{ID},\mathrm{SD}\}}\!\bigl(\beta_d+\beta_{g,d}\,G\bigr)\mathbf 1_{d}
      + \sum_{d\in\{\mathrm{FD},\mathrm{ID},\mathrm{SD}\}}\!\beta_{s,d}\,s\,\mathbf 1_{d},
\]
with $G$ the world 2100 GDP, $s\in\{0,1\}$ the structural-change (expanded public services) dummy, $d$ the decarbonisation pace and $\mathbf 1_d$ its indicator (fast decarbonisation \scenario{FD} is the reference level). The emissions effect of expanded public services is thus already in the base, through the term $\beta_{s,d}\,s$, which the interaction makes decarbonisation-specific.

On top of the base, two additive corrections capture the food and travel choices --- $-0.24$\,°C for cutting beef and $-0.155$\,°C for cutting flights. They are applied relative to the canonical pairing in which structural change ($s=1$) comes with both beef and flights reduced: if the respondent's beef or flight choice departs from that pairing, the corresponding $\pm0.24$ or $\pm0.155$\,°C is added back.

\paragraph{Calibrating the two corrections.} The beef value is the diet share of \citeauthor{Chancel2026}'s structural-change effect: their \scenario{1}$\to$\scenario{2} transition lowers 2100 warming by $0.44$\,°C, of which about $55\%$ is attributable to dietary change rather than to the accompanying reforestation --- combining the EAT-Lancet land-use saving of Hayek et al.\ (2021) with \citeauthor{Chancel2026}'s agricultural term --- giving $0.55\times0.44=\mathbf{0.24}$\,°C for the $\sim$60\,\% beef cut. The direct food-warming estimate of Ivanovich et al.\ (2023) is smaller ($\approx 0.16$\,°C at a $9.4$\,bn population) because it omits the carbon opportunity cost of the freed grazing land.

The flight value starts from $0.29$\,°C of 2100 aviation warming under high-traffic, fast-efficiency assumptions (Smith et al.\ 2026). Halving flights scales this by the scenario's projected air traffic --- a product of an income factor ($\approx 0.05\times$ GDP per capita in k\EUR), a transport-share factor and a decarbonisation factor, relative to the reference $1.14$ of the highest-traffic case. Since that would make the respondent's flight choice depend on the surrounding scenario, the correction is instead fixed at one benchmark (\scenario{BI1\_ID}: inequality, no structural change, intermediate decarbonisation), giving $\mathbf{0.155}$\,°C.

% \note{%
%   \scenario{B90k} (40\,h) is not in \citet{Chancel2026}'s emission output, so its temperature is
%   predicted rather than observed. Its 2100 GDP sits between convergence and the 44\,h class but on a
%   distinct productivity path, so this prediction is the least certain of the lot.
% }

% ==============================================================================
\section{What the browser computes at display time}
\label{sec:browser}
% ==============================================================================

Given the precomputed tables, the web page does only light arithmetic.

\paragraph{The respondent's own 2035 income.} It converts the household's monthly cash income to an annual individual income $y$ (halving it for couples) and locates its g-percentile $g^\star$ in the 2025 cash distribution \col{IT25} by linear interpolation: writing $(g_i,\col{IT25}_i)_{i=1}^{127}$ for the g-percentile grid, for the bracket with $\col{IT25}_i \le y < \col{IT25}_{i+1}$,
\[
  g^\star \;=\; g_i + \frac{y-\col{IT25}_i}{\col{IT25}_{i+1}-\col{IT25}_i}\,(g_{i+1}-g_i).
\]
The displayed income is the same interpolation of the chosen scenario's 2035 column at $g^\star$. The search runs on the 127 g-percentiles, not coarse percentiles --- 28 of them cover the top 1\,\%, where incomes vary most --- and \col{IT25} increases across all 127 brackets, so the lookup is unambiguous (an income below the first or above the last bracket clamps to it). The chosen column is assembled on the fly from the exported base columns by applying the same hours, scope, decarbonisation and public-services factors described in Section~\ref{sec:factors}, exactly reproducing the precomputed values.

\paragraph{The 2100 bar charts.} A ``national incomes'' chart shows Italy's 2100 distribution in 21 brackets (the top 1\%, the next 4\%, and the rest divided in brackets with 5\% width); a ``global incomes'' chart shows the world income. Both are read from the 2100 table for the scenario's scope and scaled by the hours multiplier of Table~\ref{tab:coefs2100}. %All values are the cash-equivalent magnitudes (already multiplied by the 0.69 scalar).

% \note{%
%   The top bracket ($p_{99}$--$p_{100}$) can be far above the median, especially in the inequality
%   scenarios. When drawing these charts, consider truncating or log-scaling the vertical axis so the
%   top bracket does not dominate the picture.
% }

% ==============================================================================
% \section{Consistency checks and known limitations}
% \label{sec:checks}
% % ==============================================================================

% \begin{enumerate}
%   \item \textbf{National-only vs inequality.} Both sit at the same no-global average, but differ in
%     shape: national-only carries the compressed, redistributed shape while inequality carries the
%     current one. The gap between them, percentile by percentile, is precisely the effect of national
%     redistribution alone.
%   \item \textbf{Global-only vs inequality.} Global-only keeps the inequality (no-national-redistribution)
%     shape but reallocates across countries the way the global transfer does worldwide: a given income
%     faces the same rate everywhere, a subsidy at the bottom and a tax at the top.
%   \item \textbf{Italy converges to the world under \scenario{SC}.} Because convergence brings
%     countries to roughly the same per-capita income by 2100, Italy's and the world's 2100 convergence
%     distributions have similar averages --- which they do in the data.
%   \item \textbf{No structural change for inequality.} \citet{Chancel2026} do not model sectoral
%     change in the inequality scenario, so an inequality society with expanded public services is
%     given the temperature of one without the associated emissions benefit.
%   \item \textbf{The \scenario{G}/\scenario{C} tension.} As noted in Section~\ref{sec:scopes}, the
%     independent construction of the inequality and convergence bases means the four scopes do not
%     satisfy an exact accounting identity; we accept the occasional small anomaly rather than risk
%     negative incomes.
% \end{enumerate}

% ==============================================================================
% \section{The two dataset variants}
% \label{sec:variants}
% % ==============================================================================

% The script writes the 2035 distributions twice. In the main file, the B (working-hours) scenarios grow between 2025 and 2035 at a flat 1.5\,\%/year. In the alternative file, every B column is rescaled to use each country's \emph{own} modelled 2025--2035 productivity growth instead; since that growth factor cancels convergence's productivity term, the 36\,h B scenario then coincides exactly with convergence. All non-B columns are identical between the two files. \file{conjoint\_few} and \file{conjoint\_any} read the flat-growth file; \file{conjoint\_any2} reads the country-growth file.

% ==============================================================================
\section{Output files}
\label{sec:files}
% ==============================================================================

\begin{description}
  \item[\file{ineq\_IT\_2035.csv} (and \file{\_full}).] 127 g-percentiles $\times$ the 2035 cash-income columns. The plain file holds the 2025 reference \col{IT25}, the net global transfer \col{cash\_GR35} and the four hours classes $\times$ four scopes the engine reads (the self-contained survey copy drops \col{cash\_GR35}, which the engine does not use); the \file{\_full} version additionally carries the construction columns (\col{IT35}, \scenario{SC} and the food/decarbonisation variants) and every hours-class $\times$ scope combination.
  \item[\file{ineq2\_IT\_2035.csv} (and \file{\_full}).] The same, with the B scenarios using each country's own 2025--2035 productivity growth.
  \item[\file{ineq\_2100.csv} (and \file{\_full}).] 21 brackets $\times$ the 2100 income columns for Italy and the world, in 2025-cash-equivalent units.
  \item[\file{temp100.csv}.] One 2100 temperature per scenario (the \citet{Chancel2026} scenarios  plus the regression-filled ones), with the GDP, structural-change and decarbonisation labels used
    to build it.
  \item[\file{conjoint\_constants.csv}.] The scalars the browser needs: the cash-to-income ratio, the public-services tax rate, the decarbonisation factors, the temperature-adjustment coefficients, GDP per capita by hours and redistribution, the 2100 populations, the regression coefficients, and the average income of every scenario in both distribution files.
\end{description}

% ==============================================================================
\section{Reproducibility}
\label{sec:repro}
% ==============================================================================

All of the above is produced by \file{code\_simulator/questionnaire.R}, which reads only the raw \citet{Bothe2026} and \citet{Chancel2026} files (plus the pre-computed Chancel temperatures) and writes the CSVs listed in Section~\ref{sec:files}, exporting the survey subset into \file{IT\_survey/data/}. The three conjoint tasks then load these files in the browser via file{IT\_survey/scenarios.js} and require only a static web server.

% ==============================================================================
% Bibliography
% ==============================================================================
\bibliographystyle{plainnat}
\begin{thebibliography}{9}

\bibitem[Bothe et al.(2026)]{Bothe2026}
Bothe, O., Chancel, L., and co-authors (2026).
\newblock Replication package for Global Justice Report Working Paper 2026-11.
\newblock \emph{WID.world} Working Paper.
\newblock URL: \url{https://wid.world/document/replication-package-working-paper-2026-11/}

\bibitem[Chancel et al.(2026)]{Chancel2026}
Chancel, L., and co-authors (2026).
\newblock Global Justice Report 2026.
\newblock \emph{World Inequality Lab}.
\newblock Data: \url{https://wseed.world/data.html}

\bibitem[Fisher-Gethin and Chancel(2023)]{FisherGethin2023}
Fisher-Gethin, A. and Chancel, L. (2023).
\newblock Distributional National Accounts: Methods and Estimates for 50 Countries.
\newblock \emph{Quarterly Journal of Economics} (forthcoming).

\bibitem[Garbinti et al.(2021)]{Garbinti2021}
Garbinti, B., Goupille-Lebret, J. and Piketty, T. (2021).
\newblock Accounting for Wealth Inequality Dynamics: Methods, Estimates and
Simulations for France.
\newblock \emph{Journal of the European Economic Association} 19(1), 620--663.

\end{thebibliography}

\end{document}
